\bigskip
\bigskip
\subsubsection*{Խնդիրներ և հարցեր թեմա 9-ի վերաբերյալ}

\begin{enumerate}[label=\thesection.\arabic*.]

% 9.1
\item Դիցուք ունենք կամայական $f: X \to Y$ անընդհատ արտապատկերում։ Ճի՞շտ է արդյոք պնդումը․ ամեն մի $x_0 \in X$ կետի ցանկացած $U$ շրջակայքի $f(U)$ կերպարը շրջակայք է $f(x_0) \in Y$ կետի համար։

\begin{hint}
Դիտարկեք որևէ $f: X \to Y$ հաստատուն արտապատկերում, որ\-տեղ $Y$-ը օժտված է անտիդիսկրետ տոպոլոգիայով և պարունակում է մեկից ավելի թվով տարրեր։
\end{hint}

% 9.2
\item Դիցուք $X$ տոպոլոգիական տարածությունը օժտված է հետևյալ հատկությամբ․ ցանկացած $Y$ տոպոլոգիական տարածության և ցանկացած $f: X \to Y$ արտա\-պատկերման դեպքում $f$-ը անընդհատ է։ Ապացուցեք, որ $X$-ի տոպոլո\-գիան դիսկրետն է։

\begin{hint}
Որպես $Y$ վերցրեք $X$ բազմությունը դիսկրետ տոպոլոգիայով։
\end{hint}

% 9.3
\item Դիցուք $Y$ տոպոլոգիական տարածությունը օժտված է հետևյալ հատկությամբ․ ցանկացած $X$ տոպոլոգիական տարածության և ցանկացած $f: X \to Y$ արտա\-պատկերման դեպքում $f$-ը անընդհատ է։ Ապացուցեք, որ $Y$-ի տոպոլո\-գիան անտիդիսկրետն է։

\begin{hint}
Որպես $X$ վերցրեք $Y$ բազմությունը անտիդիսկրետ տոպոլոգիա\-յով։
\end{hint}

% 9.4
\item Դիցուք $f: X \to Y$ արտապատկերումն այնպիսին է, որ $X$-ում բաց ցանկացած ենթաբազմության կերպարը բաց ենթաբազմություն է $Y$-ում։ Ապացուցեք, որ այդպիսի արտապատկերումը կարող է անընդհատ չլինել։

\begin{hint}
Դիտարկեք որևէ $X$ բազմություն $\tau$ և $\sigma$ տոպոլոգիաներով, որտեղ $\tau \subset \sigma$, իսկ որպես $f$ վերցրեք $\nuynakan_X$ նույնական արտապատկերումը։
\end{hint}

% 9.5
\item Բերեք $X$, $Y$ տոպոլոգիական տարածությունների և $f: X \to Y$ բիյեկտիվ ան\-ընդհատ արտապատկերման օրինակ, որ $f^{-1}$ հակադարձ արտապատկերումը չլինի անընդհատ։

\begin{hint}
Օգտվեք նախորդ խնդրի ցուցումից։
\end{hint}

% 9.6
\item Դիցուք $\R^2$ էվկլիդեսյան հարթության ինչ-որ $f: \R^2 \to \R^2$ արտապատ\-կե\-րում բավարարում է $\rho\left(f(x_1), f(x_2)\right) = r \cdot \rho(x_1, x_2)$ պայմանին $\forall \ x_1, x_2 \in \R^2$ կետերի դեպքում, որտեղ $r$-ը սևեռված իրական դրական թիվ է։ Ապա\-ցու\-ցեք, որ $f$-ը անընդհատ արտապատկերում է։

% 9.7
\item Ունենք $\R$ թվային ուղղի $g: \R \to \R$ անընդհատ արտապատկերում։ Հայտնի է, որ $g(x)$-ը ամբողջ թիվ է ցանկացած $x \in \R$ կետի դեպքում։ Ապացուցեք, որ $g$-ն հաստատուն արտապատկերում է։

% 9.8
\item Ունենք $f: X \to \R$ անընդհատ արտապատկերում, որտեղ $X$-ը կամայական տոպոլոգիական տարածություն է, և $a \in \R$ սևեռված թիվ։ Ապացուցեք, որ $g: X \to \R$ արտապատկերումը՝ սահմանված $g(x) = a f(x), \ x \in X$ բանաձևով, նույնպես անընդհատ է։

\begin{hint}
Ներկայացրեք $g$-ն որպես $f$ և $h: \mathbb{R} \to \mathbb{R}, \ h(r) = ar, \ r \in \mathbb{R}$ արտապատկերումների համադրույթ:
\end{hint}

% 9.9
\item Ունենք $X$ բազմության որևէ երկու $\tau$ և $\sigma$ տոպոլոգիա։ Ապացուցեք․ $X \to X$ նույնական արտապատկերումը $(X, \tau)$ և $(X, \sigma)$ տարածությունների հոմեոմոր\-ֆիզմ է այն և միայն այն դեպքում, երբ $\tau$-ն և $\sigma$-ն նույնն են։

% 9.10
\item Դիցուք $f: X \to Y$ արտապատկերումը $(X, \tau)$ և $(Y, \sigma)$ տոպոլոգիական տա\-րա\-ծությունների հոմեոմորֆիզմ է։ Ապացուցեք․ եթե որևէ $U \subset X$ ենթաբազմու\-թյուն շրջակայք է որևէ $x_0 \in X$ կետի համար, ապա $f(U) \subset Y$ ենթաբազմու\-թյունը շրջակայք է $f(x_0)$ կետի համար։

% 9.11
\item Ճի՞շտ է արդյոք պնդումը․ վերջավոր թվով տարրերից կազմված տոպոլոգիա\-կան տարածությունների դեպքում տարրերի քանակությունը տոպոլոգիական ինվարիանտ է։

% 9.12
\item Ընտրեք անջատելիության $\textrm{T}_0$, $\textrm{T}_1$, $\textrm{T}_2$ աքսիոմներից որևէ մեկը և ապացուցեք, որ այն տոպոլոգիական հատկություն է։

% 9.13
\item Բերեք տոպոլոգիական տարածությունների որևէ $f: X \to Y$ փակ, բայց ոչ բաց արտապատկերման օրինակ։

\begin{hint}
Դիտարկեք $\R $ թվային ուղիղը և $\R ^2$ կոորդինատային հարթու\-թյու\-նը սովորական (էվկլիդյան) մետրիկային տոպոլոգիաներով։ Ապացուցեք, որ $f: \R  \to \R ^2, \ f(x)=(x, 0), \ x \in \R $ արտապատկերումը փակ, բայց ոչ բաց արտապատկերում է։
\end{hint}

% 9.14
\item Ունենք $f: X \to Y$ և $g: Y \to Z$ արտապատկերումներ, ընդ որում ${g \circ f: X \to Z}$ համադրույթը անընդհատ է։ Ապացուցեք․ եթե $f$-ը սյուրյեկտիվ բաց (կամ փակ) արտապատկերում է, ապա $g$-ն անընդհատ արտապատկերում է։

% 9.15
\item Ճի՞շտ է արդյոք, որ թվային ուղղի $\R  \to \R $ նույնական արտապատկերումը որո\-շում է տոպոլոգիական տարածությունների $(\R, \mapsfrom) \to (\R, \mapsto)$ հոմեոմոր\-ֆիզմ, որտեղ $\mapsfrom$ և $\mapsto$ սիմվոլները թվային ուղղի ձախից կիսաբաց և աջից կի\-սա\-բաց ինտերվալների տոպոլոգիաներն են։

% 9.16
\item Ապացուցեք․ $(\R, \mapsto)$ և $(\R, \mapsfrom)$ տոպոլոգիական տարածությունները հոմեոմորֆ են։

\begin{hint}
Դիտարկեք $f: \R \to \R, \ f(x) = -x$ արտապատկերումը և օգտվեք \hyperref[ch9thm10]{թեորեմ 10}-ից։
\end{hint}

% 9.17
\item Ճի՞շտ է արդյոք, որ $(\R, \mapsto)$ և $(\R, \text{սովոր.})$ տարածությունները հոմեոմորֆ են։

\begin{hint}
Ապացուցեք, որ $[a, b)$ տեսքի ենթաբազմությունները միաժա\-մա\-նակ բաց և փակ ենթաբազմություններ են $(\R, \mapsto)$ տարածությունում, և օգտը- \linebreak վեք \hyperref[ch9thm10]{թեորեմ 10}-ից։
\end{hint}


% 9.18
\item Ունենք որևէ $f$, $g: X \to Y$ անընդհատ արտապատկերումներ, որտեղ $Y$-ը հաուսդորֆյան տարածություն է։ Դիտարկենք $F \subset X$ ենթաբազմությունը կազ\-մված այն բոլոր $x$ կետերից, որ $f(x) = g(x)$։ Ապացուցեք, որ $F$-ը փակ ենթա\-բազ\-մություն է $X$-ում։

\begin{hint}
Ցույց տվեք, որ $F$-ը պարունակում է իր բոլոր հպման կետերը։
\end{hint}

% 9.19
\item Դիցուք $f$, $g: X \to Y$ արտապատկերումները, $Y$ տարածությունը և ${F \subset X}$ ենթաբազմությունը այնպիսին են, ինչպես նախորդ խնդրում։ Ապացուցեք․ եթե $F$-ը նաև ամենուրեք խիտ է $X$-ում, ապա $f$ և $g$ արտապատկերումնե\-րը համընկնում են $X$-ի բոլոր կետերում։

\end{enumerate}
